\clearpage
\section{Concrete parametricity and the final segment of Figure 22}
\label{sec:proposition-final-segment}
The second milestone instantiates parametricity with the actual private
closure used by Figure~22 and composes it with the preceding section's
normalization. Its endpoint is the unchanged right-hand term translation,
in arbitrary finite environments and at every adequate common allowance.
It also allows any certified one-input private term sharing tree, so the
remaining source normalization need not first make that private tree canonical.

\paragraph{The remaining boundary.}
This theorem starts at the \emph{normalized private closure}, not at the
original allocating left-hand translation. Section~\ref{sec:proposition-opening}
already reaches the latter's actual opened graph. The structural bridge
between that graph and the normalized private closure remains to be proved.
No theorem here assumes that missing equivalence, and none is named or
presented as the completed Proposition~12.2.

\subsection{The actual single-closure family instance}
A family with one large member and zero small members contains exactly the
retained large member. The singleton-family isomorphism uses its original
occurrence, whose input and output maps fix the ordered boundary. Every
edge of the whole family lies in that occurrence: the other branch of the
one-copy tensor and the entire zero-copy family have empty edge carriers.
The existing occurrence theorem therefore supplies a whole-graph recursive
isomorphism. Its applications give both the single large left family and
the single large right family. No empty component is treated as a new
identity wire or operational step.
\lean{Copying/PropositionParametricityInstance}{13}{30}{ParametricInstance.singletonFamily_iso}
\lean{Copying/PropositionParametricityInstance}{32}{35}{ParametricInstance.leftFamily_iso}
\lean{Copying/PropositionParametricityInstance}{37}{43}{ParametricInstance.rightFamily_iso}

\termdef{normalized-private-closure}{Normalized private closure}.
Let $L_d$ be the large schematic read-after-callback abstraction and let
$C$ be any actual finite acyclic one-input term contraction tree. The
normalized private closure is
\[
 Q_{d,C}=(L_d;C);\mathsf{privateCell}_{d+2}.
\]
The private cell is the original instance followed by its full canonical
atom tree, store edge and numeral one. Its interface has one term input and
no outputs, so the private reference does not escape through a free
resource output. The canonical specialization chooses $C=C_1$; the more
general constructor retains $C$ without further normalization.
\lean{Copying/PropositionParametricityInstance}{44}{49}{ParametricInstance.privateClosureWith}
\lean{Copying/PropositionParametricityInstance}{50}{53}{ParametricInstance.privateClosure}

The arbitrary-tree pair proof first uses the single-left-family isomorphism
under the unchanged tree and store suffix. It then applies the literal
parametricity family constructor with $m_1=1$ and $m_2=0$, and uses the
inverse single-right-family isomorphism. This constructor asserts a
syntactic template instance. Its semantic soundness is supplied by the
already completed parametricity proof, including copying and garbage
collection dependencies. The canonical pair specializes this construction.
\lean{Copying/PropositionParametricityInstance}{54}{60}{ParametricInstance.pairWith}
\lean{Copying/PropositionParametricityInstance}{61}{62}{ParametricInstance.pair}

\subsection{Retaining both quantitative directions}
The private-program record uses the actual pair's proved focus and hole
freedom. Applying the original pair-contextual theorem yields directed
improvement from the canonical private closure to the normalized constant
closure. For the converse, place the same literal pair in each closing
context and apply the existing reverse-overhead filling theorem. Both
comparisons use the complete original machine and observation class.
Successful termination equivalence follows after forgetting the costs.
\lean{Copying/PropositionParametricityInstance}{64}{68}{ParametricInstance.privateProgram}
\lean{Copying/PropositionParametricityInstance}{69}{74}{ParametricInstance.improvement}
\lean{Copying/PropositionParametricityInstance}{75}{81}{ParametricInstance.reverse_overhead}
\lean{Copying/PropositionParametricityInstance}{83}{88}{ParametricInstance.equivalent}

Compose this concrete equivalence with the reverse of the right-hand
normalization theorem. The result runs from $Q_{d,C_1}$ to the actual
padded closed $P_R$. Directed improvement also survives the composition:
the normalization contributes equal costs, which satisfy the required
inequality reflexively. The separately proved reverse-overhead statement
starts with the exact normalization in the opposite direction and then
uses reverse parametricity. Thus no exact-cost claim is made for private-state
elimination itself.
\lean{Copying/PropositionParametricityInstance}{89}{95}{ParametricInstance.private_original}
\lean{Copying/PropositionParametricityInstance}{96}{100}{ParametricInstance.private_original_improvement}
\lean{Copying/PropositionFinalSegment}{44}{51}{ParametricInstance.private_original_reverse_overhead}

\subsection{The full final segment in arbitrary environments}
\termdef{figure22-final-segment}{Figure 22 final segment}.
The ambient private-program record adds only the actual unused term and
atom resources to the normalized private closure. Lift the closed concrete
parametricity equivalence through the certified sole-hole ambient context,
then compose with the reverse of the original-environment right normalization.
This proves the final segment for every $d,k,h$:
\[
 \mathsf{ambient}_{k,h}(Q_{d,C_1})
 \simeq R_d^{\mathrm{ambient}(k,h)}
 \simeq (P_R\text{ in }k,h)^{\mathrm{pad}(d+2)}.
\]
The endpoint is the literal unchanged translation; there is no hypothesis
postulating translation correctness, robustness or machine determinism.
\lean{Copying/PropositionFinalSegment}{9}{13}{ParametricInstance.privateEnvironment}
\lean{Copying/PropositionFinalSegment}{14}{18}{ParametricInstance.private_original_environment}

The named theorem uses the lengths of the original distinct variable and
atom lists. For any externally chosen $N\geq2$, the allowance theorem
chooses $d=N-2$, proves $d+2=N$, and gives the same equivalence at that
allowance. The observation corollary unfolds contextual equivalence into
successful-termination equivalence for every binding-free closing context.
The depth-three specialization has exactly the \texttt{translatedProgram}
endpoint needed to compare with the original left program's allowance.
\lean{Copying/PropositionFinalSegment}{20}{27}{ParametricInstance.private_original_named}
\lean{Copying/PropositionFinalSegment}{28}{37}{ParametricInstance.private_original_at}
\lean{Copying/PropositionFinalSegment}{38}{43}{ParametricInstance.private_original_termination}
\lean{Copying/PropositionFinalSegment}{52}{57}{ParametricInstance.common_depth_three}

For an arbitrary certified private tree $C$, package $Q_{d,C}$ with the
focus and hole invariants of its actual pair. Apply the same parametricity
and ambient-context theorems and compose with the same right normalization.
The general environment theorem therefore retains the entire supplied tree;
its hypothesis is the concrete contraction-tree data and acyclicity proof,
not an assumed equivalence of that tree with a canonical presentation.
\lean{Copying/PropositionFinalSegment}{58}{59}{ParametricInstance.privateProgramWith}
\lean{Copying/PropositionFinalSegment}{61}{69}{ParametricInstance.private_original_with_tree}

\subsection{Concrete regression statements}
The application counter inspects only actual outer labels of an
allowance-one body. The original translated right body, the mixed body and
the normalized schematic body each contain two application edges. Hence
the callback computation has not been removed during structural normalization.
The corresponding total edge inventories are sixteen, fifteen and seven;
the isolated sharing component has ten edges. These checked counts accompany
the actual incidence and template proofs; counts alone are never used to
justify an isomorphism or a rewrite.
\lean{Copying/PropositionFinalSegmentChecks}{10}{13}{FinalSegmentChecks.applications}
\lean{Copying/PropositionFinalSegmentChecks}{15}{20}{FinalSegmentChecks.callback_computation_retained}
\lean{Copying/PropositionFinalSegmentChecks}{21}{28}{FinalSegmentChecks.sharing_edge_inventory}

The depth-three canonical private closure has nine outer edges, while its
constant closure has two. The genuine parametricity step accounts for this
state elimination. A mixed-environment regression uses two variables and
three atoms and supplies the ten-edge tree extracted from the original
right translation as a \emph{noncanonical} private tree. It still reaches
the unchanged translated right program at allowance three.
\lean{Copying/PropositionFinalSegmentChecks}{29}{34}{FinalSegmentChecks.private_edge_inventory}
\lean{Copying/PropositionFinalSegmentChecks}{35}{41}{FinalSegmentChecks.noncanonical_mixed_environment}

Finally, for every exact step count $n$ and every binding-free observation,
the right normalization has a successful $n$-step execution on one side
if and only if it has one on the other. In each direction the exact-cost
refinement produces a matching run whose cost equality identifies its
length with $n$. This quantification includes repeated closure use and
observations with divergent or stuck callbacks; it is not a test assuming
that a selected callback terminates.
\lean{Copying/PropositionFinalSegmentChecks}{42}{56}{FinalSegmentChecks.exact_observation_count}

\subsection{Dependency map and remaining work}
With $O$ denoting the previously proved actual opened left graph, the current
chain is
\[
 P_L\ \underbrace{\simeq}_{\text{proved allocation/beta}}\ O
 \ \underbrace{\dashrightarrow}_{\text{open structural bridge}}\ Q_{d,C}
 \ \underbrace{\simeq}_{\text{proved parametricity}}\ R_d
 \ \underbrace{\simeq}_{\text{proved normalization}}\ P_R.
\]
The last two links now have concrete endpoints, arbitrary finite-environment
transport and all adequate common allowances. Their dependencies are the
original Idem/BPullW weakening certificates, actual finite graph
isomorphisms, the proved sharing-tree theorem, the proved full parametricity
family and the certified ambient context. No application-level semantic
assumption remains in these links.

The open source bridge must normalize the opened translated body and its
retained sharing and boxes to the schematic read-after-callback closure,
while preserving the actual private cell. The original-left prefix and that
bridge must be carried to the required common allowances before composing
the complete public theorem. The correspondence, substitution and any
required stability obligations in Appendix~\ref{app:remaining-work} also
remain. Proposition~12.2 is not yet proved. Completion remains estimated at
70--80\% of mathematical effort, with low confidence: these are two completed
concrete links, but the remaining bridge and correspondence work can still
require substantial mathematics.
